%!TEX root = parita-msc.tex

\chapter{Introduction}
\label{chap:intro}

\begin{doublespace}

Tim Berners-Lee is the father of the \ac{www}. 
He wrote 
\emph{Information Management: A Proposal}\footnote{{\tt https://www.w3.org/History/1989/proposal.html}, dated March 1989, May 1990} 
in 1989 and 1990 to present the need for a 
``global hypertext system'' at the \ac{cern}. 
In his conclusion to that document, he wrote:
\blockquote
{``we should work toward a universal linked information system, 
in which generality and portability are more important than fancy graphics techniques and complex extra facilities. 
The aim would be to allow a place to be found for any information or reference which one felt was important and a way of finding it afterward. 
The result should be sufficiently attractive to use, that is, 
the information contained would grow past a critical threshold so that the usefulness of the scheme would, in turn, 
encourage its increased use.''}

In the more than 30 years that have followed his initial proposal,
a great deal of progress towards Berners-Lee's vision has been realized. 
The original 
``web of documents'' is becoming the 
``web of data'' or the
``semantic web''. 

The gold standard for data on the semantic web is
5-star\footnote{\tt https://5stardata.info/en/} \ac{lod}.
Table~\ref{tab:fivestars} illustrates the progression of stars.
All resource triples (subjects, predicates, and objects) 
are named with \ac{https} \acp{uri} that can be looked-up on the web
to provide useful information with the \ac{rdf} family of standards.
This data is accompanied with an open license for access.


\begin{table}[htp]
\caption{The meaning of the stars in Berners-Lee's 5-star system. If the data is on the web without an open license, one can think of 5-star linked data (that is not open).}
\begin{center}
\begin{tabular}{|r|p{10cm}|}
\hline
Stars & Data Qualities \\ \hline
$\largestar$ & available on the Web (whatever format) under an open license \\ \hline
$\largestar\largestar$ & available as structured data (e.g., Excel instead of image scan of a table) \\ \hline 
$\largestar\largestar\largestar$ & available in a non-proprietary open format (e.g., CSV instead of Excel)\\ \hline
$\largestar\largestar\largestar\largestar$ & uses URIs to denote things, so that people can point at the contents of the data\\ \hline
$\largestar\largestar\largestar\largestar\largestar$ & links to other data to provide context\\ \hline
\end{tabular}
\end{center}
\label{tab:fivestars}
\end{table}%

Is this next web of data sufficiently attractive to use, and by whom?
Have we reached the point that Berners-Lee envisioned, 
when the quantity of information stored in the web of data, 
using \ac{rdf},
has grown past a critical threshold and the usefulness of \ac{rdf} and related technologies 
encourages their increased use?
Conservatively, there are now several billion triples on the web.
Does this wealth of data compel would-be publishers to put their own
\ac{rdf} triples onto the web? Has data consumption and production become,
for institutions and citizens, not merely accessible but welcoming?
This thesis will examine these questions in the context of an example.

\section{Motivation}
\label{sec:motivation}

In his teaching practise, 
my thesis supervisor\footnote{Daryl H. Hepting, Professor of Computer Science, University of Regina, Regina, Canada.} collects feedback from students in his classes each semester and makes the responses available on his 
website\footnote{\tt http://www2.cs.uregina.ca/$\sim$hepting/teaching/evaluation.html}. 
Hepting's feedback 
instrument\footnote{\tt http://www2.cs.uregina.ca/$\sim$hepting/assets/teaching/pdf/feedback-instrument.pdf}
(as shown in Figure~\ref{feedback-form} on page~\pageref{feedback-form})
has 9 Likert~\cite{johns2010likert} items and 3 open-ended questions.
Typically, student responses are collected on paper 
and then entered manually into computer-based spreadsheets. 
Data collection was changed and interrupted during the COVID-19 pandemic. 

Although it is valuable to have this feedback data available on the web
in any form (earning 1, 2 or 3 stars, see Table~\ref{tab:fivestars}), 
it is not published as \ac{lod}. 
The motivation for the research presented here is to examine the process 
by which existing data, in spreadsheet form, can be published as \ac{lod} 
and to consider whether that process is ``sufficiently attractive to use''.

\section{Contribution}
\label{sec:contribution}

Within the \ac{lod} community, 
most research efforts deal with users who are domain experts.
This thesis is focused on the usability of tools and frameworks 
for the development, description and reuse 
of \ac{lod} by citizens.
 
The research presented here is a practical exploration of the process of
taking spreadsheet data and publishing it on the web as 5-star \ac{lod}.
More than converting the data into \ac{rdf}, 
it is about conveying as much as possible of the semantics of the data.
This thesis examines the linked data lifecycle from the perspective of usability
for citizens.

In order to express the semantics of Likert items on Hepting's feedback instrument, 
a thorough search for vocabularies that might be reused was made of vocabulary repositories~\cite{cox2021ten} that included:
Linked Open Vocabularies (LOV)\footnote{\tt https://lov.linkeddata.es/dataset/lov}, 
Research Vocabularies Australia (RVA)\footnote{\tt https://vocabs.ardc.edu.au/}, 
BioPortal\footnote{\tt https://bioportal.org}, 
Agroportal\footnote{\tt http://aims.fao.org/agroportal}, 
Ecoportal\footnote{\tt http://ecoportal.lifewatchitaly.eu/ontologies}, 
Zenodo\footnote{\tt https://zenodo.org}, 
FAIRsharing\footnote{\tt https://fairsharing.org/} and 
Dryad\footnote{\tt https://datadryad.org/stash}.
No matches were found, so a
preliminary vocabulary to express the needed Likert semantics was developed and
it is used to represent Hepting's feedback response data. 
%To facilitate such a representation, it is important to have a clear idea 
%of the data model and the relationships amongst the data.

The publishing process has been aided by open-source tools,
namely Protégé~\cite{billey2019foundations} and OpenRefine~\cite{fiorelli2021lifting}.
Details of the practical application of these tools also represents an important contribution.

\section{Organization}
\label{sec:organization}

Following this introduction, 
the rest of the thesis is organized as follows. 
Chapter 2 describes related work that is pertinent to the thesis. 
In particular, 
concepts of the semantic web, linked open data, and related standards and technologies 
are presented. 
Chapter 3 describes the linked data lifecycle and some of the tools
required to describe semantics and map survey data. 
Chapter 4 details the experience of applying the linked data lifecycle to real data.
Finally, Chapter 5 presents some conclusions and opportunities for future work.

\end{doublespace}
